\section{Additional Evaluation Context}
\label{app:context}
\subsection{Evaluation Units and Pairing}
The main LIBERO comparison contains 2,000 episodes per method: four
suites, ten tasks per suite, and 50 episodes per task. Outcomes are paired
by initial state within the same evaluation framework and host. The
success-rate difference is reported as method minus reference in
percentage points. The paired confidence intervals and exact two-sided
McNemar tests use the discordant outcomes, that is, episodes in which
only one member of a pair succeeds. A test against the reference and a
direct test between two compressed methods answer different questions;
one method's significant difference from the reference does not establish
a significant difference between the two methods.

The 200-episode rANS deployment check is a separate sample, with 50
episodes per suite. Its payload measurements are associated with the
primary configuration but do not establish that all 2,000 success-rate
episodes traversed the deployed entropy-coding path. Episode counts,
payload samples, and upload timing samples therefore retain separate
denominators throughout the evaluation.

\subsection{Input Protocols and Evaluation Windows}
The RoboCasa input protocol is part of the experiment definition.
Native 224 rendering, unprocessed 256 rendering, and training-domain AV1
processing produce different inputs to the same policy. Their reference
success rates are not interchangeable. The comparison in
Table~\ref{tab:robocasa} uses the training-domain protocol; the pooled
reference and proposed-method result combines two batches, with
comparisons matched within each batch and host.

Each RoboCasa batch contains 24 tasks and 20 episodes per task. The first
batch ran on A100 and the second on A40; their reference and proposed
method counts were 291 versus 295 and 286 versus 275, respectively.
The pooled 960-episode result combines the paired comparisons within
these batches. It does not pair outcomes across hosts. In the first
batch, the reference, SPARC, and H.264 clients shared a policy service,
so matching initial states did not align the policy's random-number
sequence for every invocation.

OpenVLA-OFT results also belong to distinct evaluation windows. The
initial complete window and the partial contemporaneous rerun are
reported independently because the reference policy exhibited drift.
A codec trained for another backbone is a transfer diagnostic rather
than a policy-matched primary baseline.

\subsection{Perturbation Breakdown}
Table~\ref{tab:appendix-perturbations} gives the category counts behind
the 700-episode LIBERO-Plus result. Each category contains 100 episodes.
The proposed codec was trained on the original LIBERO observations.
Its largest observed deficits relative to the reference were nine
successes under camera-viewpoint changes and seven under background
texture changes. The table describes where the observed deficit occurs;
it does not isolate its cause. Individual category counts also do not
establish a general robustness ranking across policies or perturbation
distributions.

\begin{table}[t]
\caption{LIBERO-Plus successes by perturbation category. Each category
has 100 paired episodes. The total row covers 700 episodes per method.}
\label{tab:appendix-perturbations}
\centering\small
\setlength{\tabcolsep}{4pt}
\begin{tabularx}{\columnwidth}{Xrrr}
\toprule
Perturbation & Reference & Ours & SPARC \\
\midrule
Camera viewpoint & 85 & 76 & 76 \\
Background texture & 96 & 89 & 83 \\
Language instruction & 93 & 89 & 92 \\
Lighting & 94 & 96 & 83 \\
Object layout & 84 & 85 & 83 \\
Robot initial state & 78 & 76 & 67 \\
Sensor noise & 80 & 79 & 74 \\
\midrule
Total & 610 & 590 & 558 \\
\bottomrule
\end{tabularx}
\end{table}

\section{Measurement Scope}
\label{app:measurement}
\subsection{Payload Accounting}
The LIBERO rANS subset yielded mean payloads of 1,857, 1,599, 2,035,
and 1,787 bytes per upload for Spatial, Object, Goal, and Long,
respectively. These measurements cover both camera streams. The
reported summary is approximately 1.82\,kB per upload; its scope is
codec payload rather than complete network traffic.

For SO-101 replay, the reporting script averages the per-upload
\texttt{payload\_bytes} field after warm-up. It divides values of at
least 1,024 bytes by 1,024 for display. Table~\ref{tab:replay} labels
these values as KiB and retains bytes for the smaller SPARC payloads.
This unit convention is distinct from the simulation summaries.

\subsection{Timing and Aggregation}
The replay discards the first five of 200 uploads for each method,
leaving 195 uploads. For each timing field, the reporting script sorts
the available values and selects the element at the nearest integer
index to the requested fraction of the zero-based index range. It does
not interpolate between adjacent observations. With 195 valid values,
p50 and p95 are the 98th and 185th sorted observations, respectively.
These are empirical upload quantiles rather than confidence limits.

Cloud-total time includes cloud decoding and policy computation; for
image paths, it also includes the original visual encoder. Cloud
decoding is therefore already included in cloud-total time. The network
residual subtracts edge encoding and cloud-total time from the complete
request duration, retaining transfer, response return, and serialization
effects. It is not a separately timed one-way network delay.

End-to-end quantiles are computed from complete request measurements.
Summing stage medians would generally combine values from different
uploads and does not recover the median request duration. Likewise,
synchronized component microbenchmarks and the captured deployed path
have different timing boundaries, as described in
\S\ref{sec:implementation}.

\subsection{Physical Interpretation}
The replay experiment uses recorded observations without corresponding
physical task execution. It can characterize request processing under
the measured trace but cannot quantify recovery from an execution error,
success under moving objects, or the age of an observation when a
physical actuator applies the returned action.

\section{Use of Generative Writing Assistance}
Generative AI assisted in drafting and editing the manuscript text and
in preparing an initial architecture illustration. The experimental
results are drawn from the project's recorded summaries; the writing
process did not generate or rerun those experiments.
